\documentclass[conference]{IEEEtran}
\IEEEoverridecommandlockouts
\usepackage{cite}
\usepackage{amsmath,amssymb,amsfonts}
\usepackage{algorithmic}
\usepackage{graphicx}
\usepackage{textcomp}
\def\BibTeX{{\rm B\kern-.05em{\sc i\kern-.025em b}\kern-.08em
    T\kern-.1667em\lower.7ex\hbox{E}\kern-.125emX}}
\begin{document}

\title{Semi-Supervised DeepSDF with Eikonal Regularization and Fourier Positional Encoding}

\author{\IEEEauthorblockN{Author Name}
\IEEEauthorblockA{\textit{School of Computer Science and Engineering} \\
\textit{Nanyang Technological University}\\
Singapore \\
email@ntu.edu.sg}
}

\maketitle

\begin{abstract}
% ~150 words. IEEE: no symbols, special characters, footnotes, or math in the title or abstract.
Replace this placeholder with the abstract summarizing the research question, method, and main findings (Eikonal label efficiency, Fourier PE failure mode).
\end{abstract}

\begin{IEEEkeywords}
DeepSDF, signed distance functions, Eikonal regularization, semi-supervised learning, positional encoding, 3D reconstruction
\end{IEEEkeywords}

\section{Introduction}
% Motivation: semi-supervised SDF learning, label efficiency for 3D reconstruction.

\section{Related Work}
% DeepSDF; Eikonal / geometric regularization in neural implicits; PE (NeRF, SIREN, etc.).

\section{Method}
% Architecture, $L_{\text{sdf}}$, $L_{\text{eik}}$, $L_{\text{2nd}}$, $L_z$, training procedure, Fourier PE formulation.

\section{Experiments}
% ShapeNet setup, metrics (CD, NC), experiment matrix, training hyperparameters.

\section{Results}
% Label-efficiency curve, PE ablation, multi-seed validation, per-category analysis if applicable.

\section{Discussion}
% Why Eikonal helps; why PE fails (extrapolation / high-frequency oscillations); limitations.

\section{Conclusion}
% Practical takeaway (e.g., low labels + Eikonal); future work.

% Uncomment and place after first citation in text:
% \section*{Acknowledgment}
% This work was supported by ...

\begin{thebibliography}{00}
\bibitem{deepsdf} J. J. Park et al., ``DeepSDF: Learning continuous signed distance functions for shape representation,'' in \textit{Proc. IEEE CVPR}, 2019.
% Add remaining references (NeRF, Tancik Fourier features, course materials, etc.).
\end{thebibliography}

\end{document}
